% Part: turing-machines 
% Chapter: undecidability 
% Section: verification

\documentclass[../../../include/open-logic-section]{subfiles}

\begin{document}

\olfileid{tur}{und}{ver} 
\olsection{Mriksa Representasi}

\begin{explain}
Kanggo mriksa yen representasi kita bisa digunakake,
kita kudu mbuktekake rong bab. Kapisan, kita kudu
nuduhake yen $M$ mandheg nganggo input~$w$, mula
$!T(M, w) \lif !E(M, w)$ sah. Banjur kita kudu
nuduhake kosok baline: yen $!T(M, w) \lif !E(M, w)$
sah, mula $M$ pancen pungkasane mandheg nalika
dilakokake nganggo input~$w$.

Strategi kanggo mbuktekake rong bab iki béda banget.
Kanggo asil kapisan, kita kudu nuduhake yen
!!a{sentence} logika orde kapisan (yaiku
$!T(M, w) \lif !E(M, w)$) sah. Cara paling
gampang yaiku menehi !!a{derivation}. Nanging
bukti kita kudu lumaku kanggo kabeh $M$ lan $w$,
mula dudu mung siji !!{sentence} kang mbutuhake
!!a{derivation}, nanging tanpa wates cacahe.
Mula kang bisa kita tindakake yaiku mbuktekake
kanthi induksi yen, apa wae wujud $M$ lan~$w$
lan pira wae langkah kang dibutuhake $M$
kanggo mandheg nganggo input~$w$, mesthi ana
!!a{derivation} saka $!T(M, w) \lif !E(M, w)$.

Mesthi wae, induksi kita bakal adhedhasar cacahe
langkah kang ditindakake $M$ sadurunge tekan
konfigurasi mandheg. Ing bukti induktif iki,
kita bakal netepake yen kanggo saben langkah~$n$
nalika $M$ lumaku nganggo input~$w$,
$!T(M, w) \Entails !C(M, w, n)$, ing ngendi
$!C(M, w, n)$ njlentrehake konfigurasi $M$
sawisé lumaku nganggo input~$w$ nganti~$n$
langkah. Yen $M$ mandheg nganggo input~$w$
sawisé, upamane, $n$ langkah, $!C(M, w, n)$
bakal njlentrehake konfigurasi mandheg.
Kita uga bakal nuduhake yen $!C(M, w, n)
\Entails !E(M, w)$ saben $!C(M, w, n)$
njlentrehake konfigurasi mandheg. Mula yen
$M$ mandheg nganggo input~$w$, ana sawijining~$n$
saengga $M$ ana ing konfigurasi mandheg sawisé
$n$~langkah. Dadi $!T(M, w) \Entails !C(M, w, n)$,
ing ngendi $!C(M, w, n)$ njlentrehake konfigurasi
mandheg; lan amarga ing kahanan iki
$!C(M, w, n) \Entails !E(M, w)$, kita entuk
% OLPL-156: Sumber pungkasan iki nulis T(M,w) tanpa tandha formula, sanajan premis lan kesimpulan sadurunge nggunakake !T(M,w).
$T(M, w) \Entails !E(M, w)$, yaiku
$\Entails !T(M, w) \lif !E(M, w)$.

Strategi kanggo kosok baline béda banget.
Ing kene kita nganggep yen
$\Entails !T(M, w) \lif !E(M, w)$
lan kudu mbuktekake yen $M$ mandheg
nganggo input~$w$. Saka hipotesis iki
kita entuk $!T(M, w) \Entails !E(M, w)$,
yaiku $!E(M, w)$ bener ing saben
!!{structure} kang ndadekake $!T(M, w)$
bener. Mula kita bakal njlentrehake
!!a{structure}~$\Struct{M}$ kang ndadekake
$!T(M, w)$ bener: domaine yaiku $\Nat$,
dene interpretasi kabeh $\Obj Q_q$ lan
$\Obj S_\sigma$ diwenehake dening
konfigurasi~$M$ sajrone lumaku nganggo
input~$w$. Contone,
% OLPL-157: Sumber ing tuladha iki nyebut T tinimbang M minangka mesin kang lumaku; identitas simbol ditahan nganti audit sumber rampung.
$\Sat{M}{\Obj Q_q(\num{m}, \num{n})}$
yen lan mung yen $T$, sawisé lumaku
nganggo input~$w$ nganti $n$ langkah,
ana ing kahanan~$q$ lan lagi maca
kothak~$m$. Amarga miturut hipotesis
$!T(M, w) \Entails !E(M, w)$,
lan amarga $\Sat{M}{!T(M, w)}$
miturut konstruksi, mula $\Sat{M}{!E(M, w)}$.
Nanging $\Sat{M}{!E(M, w)}$
yen lan mung yen ana $n \in \Domain{M} = \Nat$
saengga $M$, nalika lumaku nganggo
input~$w$, ana ing konfigurasi mandheg
sawisé~$n$ langkah.
\end{explain}

\begin{defn} 
Tetepna $!C(M, w, n)$ minangka !!{sentence}
\[ 
\Obj Q_q(\num{m}, \num{n}) \land \Obj S_{\sigma_0}(\num{0}, \num{n})
\land \dots \land \Obj S_{\sigma_k}(\num{k}, \num{n}) \land
\lforall[x][(\num{k} < x \lif \Obj S_\TMblank(x, \num{n}))]
\] 
ing ngendi $q$ yaiku kahanan $M$ ing wektu~$n$,
$M$ maca kothak~$m$ ing wektu~$n$, kothak~$i$
ngemot simbol~$\sigma_i$ ing wektu~$n$ kanggo
$0 \le i \le k$, lan $k$ yaiku kothak paling
tengen kang ora kosong ing pita ing wektu~$0$,
utawa kothak paling tengen kang wis ditekani
sirah maca-nulis sawisé $n$ langkah, endi
wae kang luwih gedhé.
\end{defn}

\begin{lem}\ollabel{lem:halt-config-implies-halt}
Yen $M$ kang lumaku nganggo input~$w$
ana ing konfigurasi mandheg sawisé $n$
langkah, mula $!C(M, w, n) \Entails !E(M, w)$.
\end{lem}

\begin{proof}
Upamane $M$ mandheg kanggo input~$w$
sawisé $n$ langkah. Ana kahanan~$q$,
kothak~$m$, lan simbol~$\sigma$ saengga:
\begin{enumerate} 
\item Sawisé $n$ langkah, $M$ ana ing
  kahanan~$q$ lan maca kothak~$m$
  kang ngemot~$\sigma$.
\item Fungsi transisi $\delta(q, \sigma)$
  ora ditegesi.
\end{enumerate}
$!C(M, w, n)$ njlentrehake konfigurasi iki
lan ngemot klausa $\Obj Q_{q}(\num{m}, \num{n})$
lan $\Obj S_{\sigma}(\num{m}, \num{n})$.
Rong klausa iki bebarengan nyebabake
$!E(M, w)$:
\[
\lexists[x][\lexists[y][(\bigvee_{\tuple{q, \sigma} \in
      X}(\Obj Q_q(x, y) \land \Obj S_\sigma(x, y)))]]
\]
% OLPL-158: Baris sumber iki ngalih saka q,sigma menyang q',sigma' tanpa ngenalake saksi anyar lan nulis S tanpa Obj; rumus beku ditahan nganti audit.
amarga $\Obj Q_{q'}(\num{m}, \num{n}) \land S_{\sigma'}(\num{m},
\num{n}) \Entails \bigvee_{\tuple{q, \sigma} \in X} (\Obj Q_q(\num{m},
\num{n}) \land \Obj S_{\sigma}(\num{m}, \num{n}))$,
amarga $\tuple{q',\sigma'} \in X$.
\end{proof}

\begin{explain}
Mula yen $M$ mandheg kanggo input $w$,
ana sawijining~$n$ saengga $!C(M, w, n)
\Entails !E(M,w)$. Saiki kita bakal
nuduhake yen kanggo saben wektu~$n$,
$!T(M, w) \Entails !C(M, w, n)$.
\end{explain}

\begin{lem}
\ollabel{lem:config}
Kanggo saben $n$, yen $M$ durung
mandheg sawisé $n$ langkah, mula
$!T(M, w) \Entails !C(M, w, n)$.
\end{lem}

\begin{proof}
Landhesan induksi: Yen $n = 0$, konjungsi-konjunksi
ing $!C(M, w, 0)$ uga ana ing $!T(M, w)$,
mula bisa didudut saka iku.

Hipotesis induksi: Yen $M$ durung mandheg
sadurunge langkah kapangka-$n$, mula
$!T(M,w) \Entails !C(M, w, n)$.
Kita kudu nuduhake yen (kajaba yen
$!C(M, w, n)$ njlentrehake konfigurasi
mandheg), $!T(M, w) \Entails
!C(M, w, n+1)$.

% OLPL-159: Langkah induksi ing sumber mung nganggep n>0, mula kasus n=0 menyang n=1 durung kacathet ing bukti iki.
Upamane $n > 0$ lan sawisé $n$ langkah,
$M$ kang diwiwiti nganggo $w$ ana ing
kahanan~$q$ lan maca kothak~$m$.
Amarga $M$ ora mandheg sawisé~$n$
langkah, mesthi ana instruksi kanthi
salah siji saka telung wujud ing ngisor
iki ing program~$M$:

\begin{enumerate} 
\item \ollabel{right} $\delta(q, \sigma) = \tuple{q', \sigma', \TMright}$

\item \ollabel{left} $\delta(q, \sigma) = \tuple{q', \sigma', \TMleft}$

\item \ollabel{stay} $\delta(q, \sigma) = \tuple{q', \sigma', \TMstay}$
\end{enumerate}

Kita bakal nliti telung kasus iki siji-siji.

\begin{enumerate} 
\item Upamane ana instruksi kanthi wujud~\olref{right}.
  Miturut \olref[rep]{defn:tm-descr}\olref[rep]{rep-right}, iki ateges
\begin{align*} 
& \lforall[x][\lforall[y][((\Obj Q_{q}(x,
  y) \land \Obj S_\sigma(x, y)) \lif {}]]\\
& \qquad (\Obj Q_{q'}(x',
  y') \land \Obj S_{\sigma'}(x, y') \land !A(x, y)))  
\intertext{iku konjungsi saka $!T(M,w)$. Saka iki
  kita entuk !!{sentence} ing ngisor iki
  (instansiasi universal, $\num{m}$ kanggo~$x$
  lan $\num{n}$ kanggo~$y$):}
& (\Obj Q_{q}(\num{m}, \num{n}) \land \Obj S_{\sigma}(\num{m},
\num{n})) \lif {}\\
& \qquad (\Obj Q_{q'}(\num{m}', \num{n}') \land
\Obj S_{\sigma'}(\num{m}, \num{n}') \land !A(\num{m}, \num{n})).
\intertext{Miturut hipotesis induksi,
  $!T(M, w) \Entails !C(M, w, n)$, yaiku}
& \Obj Q_q(\num{m}, \num{n}) \land \Obj S_{\sigma_0}(\num{0}, \num{n})
\land \dots \land \Obj S_{\sigma_k}(\num{k}, \num{n}) \land \\
& \qquad\lforall[x][(\num{k} < x \lif \Obj S_\TMblank(x, \num{n}))]\\
\intertext{Amarga sawisé $n$ langkah,
  kothak pita~$m$ ngemot~$\sigma$, konjungsi
  kang cocog yaiku~$\Obj S_\sigma(\num{m}, \num{n})$,
  mula saka iki kita entuk:}
& \Obj Q_{q}(\num{m}, \num{n}) \land \Obj S_{\sigma}(\num{m},
   \num{n})
\intertext{Saiki kita entuk}
& \Obj Q_{q'}(\num{m}', \num{n}') \land \Obj S_{\sigma'}(\num{m},
  \num{n}') \land {}\\
& \qquad\Obj S_{\sigma_0}(\num{0}, \num{n}') \land \dots \land
  \Obj S_{\sigma_k}(\num{k}, \num{n}') \land {}\\
& \qquad \lforall[x][(\num{k} < x \lif \Obj S_\TMblank(x, \num{n}'))]
\end{align*}
alesane kaya mangkene: Larik kapisan langsung
ndherek saka konsekuen implikasi sadurunge,
kanthi modus ponens. Saben konjungsi ing larik
tengah---kang ora kalebu $S_{\sigma_m}(\num{m},\num{n}')$
---ndherek saka konjungsi kang cocog ing
$!C(M, w, n)$ bebarengan karo $!A(\num{m},\num{n})$.
% OLPL-160: Sumber nyebut larik tengah ora ngemot sel m, nanging tampilan sakdurunge isih nyerat kabeh S_sigma0 nganti S_sigmak.

Yen $m < k$, $!T(M,w) \Proves \num{m} < \num {k}$
(\olref[rep]{prop:mlessk}), lan saka transitivitas~$<$,
kita entuk $\lforall[x][(\num{k} < x \lif \num{m} < x)]$.
Yen $m = k$, banjur
$\lforall[x][(\num{k} < x \lif \num{m} < x)]$
ndherek saka logika wae. Larik pungkasan banjur
ndherek saka konjungsi kang cocog ing $!C(M, w, n)$,
$\lforall[x][(\num{k} < x \lif \num{m} < x)]$,
lan $!A(\num{m},\num{n})$. Yen $m<k$,
iki wis dadi $!C(M, w, n+1)$.

Saiki upamane $m=k$. Ing kahanan iki, sawisé
$n+1$ langkah, sirah maca-nulis uga wis
ndatangi kothak~$k+1$, kang saiki dadi
kothak paling tengen kang tau ditekani.
Mula $!C(M, w, n+1)$ nduwé konjungsi anyar,
$\Obj S_\TMblank(\num{k}',\num{n}')$,
lan konjungsi pungkasan yaiku
$\lforall[x][(\num{k}' < x \lif \Obj S_\TMblank(x, \num{n}'))]$.
Kita kudu mriksa yen rong !!{sentence}s
iki uga bisa didudut.

Kita wis nduwé $\lforall[x][(\num{k} < x \lif \Obj S_\TMblank(x,
  \num{n}'))]$. Mligine, iki menehi
$\num{k} < \num{k}' \lif \Obj S_\TMblank(\num{k}', \num{n}')$.
Saka aksioma $\lforall[x][x < x']$,
kita entuk $\num{k} < \num{k}'$.
Kanthi modus ponens, kita entuk
$\Obj S_\TMblank(\num{k}',\num{n}')$.

Semono uga, amarga $!T(M,w) \Proves \num{k} < \num{k}'$,
aksioma transitivitas~$<$ menehi
$\lforall[x][(\num{k}' < x \lif \Obj
  S_\TMblank(x, \num{n}'))]$.
(Buktine iki kita tinggalake dadi latihan.)

\item Upamane ana instruksi kanthi wujud~\olref{left}.
  Miturut \olref[rep]{defn:tm-descr}\olref[rep]{rep-left},
\begin{align*} 
& \lforall[x][\lforall[y][((\Obj Q_{q}(x', y) \land \Obj
    S_{\sigma}(x', y)) \lif {}]]\\
& \qquad (\Obj Q_{q'}(x, y') \land \Obj
  S_{\sigma'}(x', y') \land !A(x, y))) \land {}\\
& \lforall[y][((\Obj Q_{q_i}(\num{0}, y) \land \Obj S_{\sigma}(\num{0},
    y)) \lif {}]\\
& \qquad (\Obj Q_{q_j}(\num{0}, y') \land \Obj S_{\sigma'}(\num{0}, y')
  \land !A(\num{0}, y)))
\intertext{iku konjungsi saka $!T(M,w)$.
  Yen $m>0$, tetepna $l = m - 1$
  (yaiku $m = l+1$). Konjungsi kapisan
  saka !!{sentence} ing ndhuwur
  menehi akibat iki:}
& (\Obj Q_{q}(\num{l}', \num{n}) \land \Obj S_{\sigma}(\num{l}', \num{n}))
\lif {} \\
& \qquad (\Obj Q_{q'}(\num{l}, \num{n}') \land \Obj S_{\sigma'}(\num{l}',
\num{n}') \land !A(\num{l}, \num{n}))
\intertext{Yen ora mangkono, tetepna
  $l = m = 0$ lan gatekna !!{sentence}
  ing ngisor iki kang ndherek saka
  konjungsi kapindho:}
& ((\Obj Q_{q_i}(\num{0}, \num{n}) \land \Obj S_{\sigma}(\num{0}, \num{n})) \lif {}\\
& \qquad   (\Obj Q_{q_j}(\num{0}, \num{n}') \land \Obj S_{\sigma'}(\num{0}, \num{n}') \land
!A(\num{0}, \num{n}))) 
\intertext{Saben ukara iki nyebabake}
&  \Obj Q_{q'}(\num{l}, \num{n}') \land \Obj S_{\sigma'}(\num{m},
  \num{n}') \land {}\\
&\qquad  \Obj S_{\sigma_0}(\num{0}, \num{n}')\land \dots \land
  \Obj S_{\sigma_k}(\num{k}, \num{n}')  \land {} \\
& \qquad  \lforall[x][(\num{k} < x
  \lif \Obj S_\TMblank(x, \num{n}'))]
\end{align*}
minangka sadurunge. (Gatekna yen ing kasus
kapisan $\num{l}' \ident \num{l+1} \ident \num{m}$,
dene ing kasus kapindho $\num{l} \ident \num{0}$.)
Nanging iki persis $!C(M, w, n+1)$.

\item Kasus \olref{stay} ditinggal dadi latihan.
\end{enumerate}
Kita wis nuduhake yen kanggo sembarang~$n$,
$!T(M, w) \Entails !C(M, w, n)$.
\end{proof}

\begin{prob}
Rampungna kasus~\olref[tur][und][ver]{stay}
ing bukti \olref[tur][und][ver]{lem:config}.
\end{prob}

\begin{prob}
Wenehana !!a{derivation} kanggo
$\Obj S_{\sigma_i}(\num{i}, \num{n}')$
saka $\Obj S_{\sigma_i}(\num{i}, \num{n})$
lan $!A(m, n)$ (kanthi nganggep $i \neq m$,
yaiku $i < m$ utawa $m < i$).
\end{prob}

\begin{prob}
Wenehana !!a{derivation} kanggo
$\lforall[x][(\num{k}' < x \lif \Obj
  S_\TMblank(x, \num{n}'))]$ saka
$\lforall[x][(\num{k} < x \lif \Obj
  S_\TMblank(x, \num{n}'))]$,
$\lforall[x][x < x']$, lan
$\lforall[x][\lforall[y][\lforall[z][ ((x < y \land y < z) \lif x <
      z)]]]$.)
\end{prob}


\begin{lem}
\ollabel{lem:valid-if-halt}
Yen $M$ mandheg nganggo input~$w$,
mula $!T(M, w) \lif !E(M, w)$ sah.
\end{lem}

\begin{proof}
Miturut \olref{lem:config}, kanggo saben
wektu~$n$, andharan~$!C(M, w, n)$
ngenani konfigurasi $M$ ing wektu~$n$
bisa didudut saka~$!T(M, w)$.
Upamane $M$ mandheg sawisé $k$ langkah.
Ing wektu iku mesin maca kothak~$m$
kanggo sawijining~$m \in \Nat$.
Mula $!C(M, w, k)$ njlentrehake
konfigurasi mandheg saka~$M$, yaiku
ngemot konjungsi $\Obj Q_q(\num{m}, \num{k})$
lan $\Obj S_\sigma(\num{m}, \num{k})$
kanthi $\delta(q,\sigma)$ ora ditegesi.
Mula, miturut \olref{lem:halt-config-implies-halt},
$!C(M, w, k) \Entails !E(M, w)$.
Nanging amarga $!T(M, w) \Entails !C(M, w, k)$,
kita entuk $!T(M, w) \Entails !E(M, w)$,
mula $!T(M, w) \lif !E(M, w)$ sah.
\end{proof}


\begin{explain} 
Kanggo ngrampungake pamriksan pratelan kita,
kita uga kudu netepake arah kosok baline:
yen $!T(M, w) \lif !E(M, w)$ sah,
mula $M$ pancen mandheg nalika diwiwiti
nganggo input~$w$.
\end{explain}

\begin{lem}
\ollabel{lem:halt-if-valid}
Yen $\Entails !T(M, w) \lif !E(M, w)$,
mula $M$ mandheg nganggo input~$w$.
\end{lem}

\begin{proof}
Gatekna $\Lang L_M$-!!{structure}~$\Struct M$
kanthi domain~$\Nat$ kang nginterpretasi
$\Obj 0$ minangka~$0$, $\prime$~minangka
fungsi penerus, $<$~minangka relasi luwih
cilik, lan predikat $\Obj Q_q$
lan~$\Obj S_\sigma$ kaya mangkene:
\begin{align*}
  \Assign{\Obj Q_q}{M} & =
\Setabs{\tuple{m, n}}{\begin{array}{ll}\text{diwiwiti nganggo~$w$, sawisé~$n$ langkah,}\\ \text{$M$ ana ing kahanan $q$
  lan maca kothak~$m$}\end{array}} \\
\Assign{\Obj S_\sigma}{M} & = \Setabs{\tuple{m, n}}{\begin{array}{ll}
\text{diwiwiti nganggo~$w$, sawisé $n$ langkah,}\\ \text{kothak~$m$ saka $M$ ngemot
  simbol~$\sigma$}\end{array}}
\end{align*}
Kanthi tembung liya, kita mbangun
!!{structure}~$\Struct{M}$ supaya
njlentrehake tumindake $M$ kang
diwiwiti nganggo input~$w$ kanthi
saben langkah. Cetha yen
$\Sat{M}{!T(M, w)}$. Yen
$\Entails !T(M, w) \lif !E(M, w)$,
mula uga $\Sat{M}{!E(M, w)}$, yaiku
\[
\Sat{M}{\lexists[x][\lexists[y][(\bigvee_{\tuple{q, \sigma} \in
      X}(\Obj Q_q(x, y) \land \Obj S_\sigma(x, y)))]]}.
\]
Amarga $\Domain{M} = \Nat$, mesthi ana
$m$, $n \in \Nat$ saengga
$\Sat{M}{\Obj Q_q(\num{m}, \num{n}) \land \Obj S_\sigma(\num{m},
\num{n})}$ kanggo sawijining~$q$
lan~$\sigma$ kang $\delta(q, \sigma)$
ora ditegesi. Miturut definisi~$\Struct M$,
iki ateges $M$ kang diwiwiti nganggo
input~$w$ sawisé~$n$ langkah ana ing
kahanan~$q$ lan maca simbol~$\sigma$,
dene fungsi transisine ora ditegesi,
yaiku $M$~wis mandheg.
\end{proof}

\end{document}
